%%%%%%%%%%%%%%%%%%%%%%%%%%%%%%%%%%%%%%%%%%%%%%%%%%%%%%%%%%%%%%%%%%%%%%%%%%%%%%%%%%
\begin{frame}[fragile]\frametitle{}
\begin{center}
{\Large Iterators}
\end{center}
\end{frame}

%%%%%%%%%%%%%%%%%%%%%%%%%%%%%%%%%%%%%%%%%%%%%%%%%%%%%%%%%%%%%%%%%%%%%%%%%%%%%%%%%%%
\begin{frame}[fragile]\frametitle{Usage of ``for'' loop}
Looping over a list:
\begin{lstlisting}
>>> for i in [1, 2, 3, 4]:
...     print(i)
\end{lstlisting}

With a string, it loops over its characters:
\begin{lstlisting}
>>> for c in "python":
...     print(c)
\end{lstlisting}
\end{frame}

%%%%%%%%%%%%%%%%%%%%%%%%%%%%%%%%%%%%%%%%%%%%%%%%%%%%%%%%%%%%%%%%%%%%%%%%%%%%%%%%%%%
\begin{frame}[fragile]\frametitle{Usage of ``for'' loop}
With a dictionary, it loops over its keys:
\begin{lstlisting}
>>> for k in {"x": 1, "y": 2}:
...     print(k)
\end{lstlisting}

With a file, it loops over lines of the file:
\begin{lstlisting}
>>> for line in open("a.txt"):
...     print(line)
\end{lstlisting}
\end{frame}

%%%%%%%%%%%%%%%%%%%%%%%%%%%%%%%%%%%%%%%%%%%%%%%%%%%%%%%%%%%%%%%%%%%%%%%%%%%%%%%%%%%
\begin{frame}[fragile]\frametitle{Iterables}
    \begin{itemize}
    \item  So there are many types of objects which can be used with a for loop. 
    \item These are called iterable objects.
	\item There are many functions which consume these iterables:
	\begin{lstlisting}
>>> ",".join(["a", "b", "c"])
'a,b,c'
>>> ",".join({"x": 1, "y": 2})
'y,x'
>>> list("python")
['p', 'y', 't', 'h', 'o', 'n']
>>> list({"x": 1, "y": 2})
['y', 'x']
\end{lstlisting}
    \end{itemize}
\end{frame}

%%%%%%%%%%%%%%%%%%%%%%%%%%%%%%%%%%%%%%%%%%%%%%%%%%%%%%%%%%%%%%%%%%%%%%%%%%%%%%%%%%%
\begin{frame}[fragile]\frametitle{Iteration}
    \begin{itemize}
    \item  The act of going over a collection
    \item  Explicit in the form of `for'
    \item  Implicit in many forms: list(), tuple(), dict(), map(), reduce(), zip()
    \item Used with `in'
    \end{itemize}
\end{frame}

%%%%%%%%%%%%%%%%%%%%%%%%%%%%%%%%%%%%%%%%%%%%%%%%%%%%%%%%%%%%%%%%%%%%%%%%%%%%%%%%%%%
\begin{frame}[fragile]\frametitle{Iteration}
    \begin{itemize}
    \item  An iterator is an object adhering to the iterator protocol : basically this means that it has a next method, which, when called, returns the next item in the sequence, 
    \item and when there's nothing to return, raises the StopIteration exception.
\item 
An iterator object allows to loop just once.
\item  It holds the state (position) of a single iteration, or from the other side, each loop over a sequence requires a single iterator object. 
\item This means that we can iterate over the same sequence more than once concurrently. 
\item Separating the iteration logic from the sequence allows us to have more than one way of iteration.
    \end{itemize}
\end{frame}

%%%%%%%%%%%%%%%%%%%%%%%%%%%%%%%%%%%%%%%%%%%%%%%%%%%%%%%%%%%%%%%%%%%%%%%%%%%%%%%%%%%
\begin{frame}[fragile]\frametitle{Iteration}
    \begin{itemize}
\item Calling the \_\_iter\_\_ method on a container to create an iterator object is the most straightforward way to get hold of an iterator. 
\item The iter function does that for us, saving a few keystrokes.
\item When used in a loop, StopIteration is swallowed and causes the loop to finish. 
\item But with explicit invocation, we can see that once the iterator is exhausted, accessing it raises an exception.
    \end{itemize}
\end{frame}

%%%%%%%%%%%%%%%%%%%%%%%%%%%%%%%%%%%%%%%%%%%%%%%%%%%%%%%%%%%%%%%%%%%%%%%%%%%%%%%%%%%
\begin{frame}[fragile]\frametitle{Iteration}
    \begin{itemize}
\item Support for iteration is pervasive in Python: all sequences and unordered containers in the standard library allow this. 
\item The concept is also stretched to other things: e.g. file objects support iteration over lines.
\item The file is an iterator itself and it's \_\_iter\_\_ method doesn't create a separate object: only a single thread of sequential access is allowed.
    \end{itemize}
\end{frame}



%%%%%%%%%%%%%%%%%%%%%%%%%%%%%%%%%%%%%%%%%%%%%%%%%%%%%%%%%%%%%%%%%%%%%%%%%%%%%%%%%%%
\begin{frame}[fragile]\frametitle{Iteration Protocol}
    \begin{itemize}
    \item Protocol of 2 methods:
    \begin{itemize}
    	\item \lstinline|__iter__()|: returns the iterator object itself
	\item \lstinline|__next__()|: returns the next value; raises \lstinline|StopIteration| when exhausted
        \end{itemize}
        \item Call via the built-in: \lstinline|next(obj)| invokes \lstinline|obj.__next__()|
        \item Explicit iterator creation: \lstinline|iter(obj)| invokes \lstinline|obj.__iter__()|
    \end{itemize}
\end{frame}

%%%%%%%%%%%%%%%%%%%%%%%%%%%%%%%%%%%%%%%%%%%%%%%%%%%%%%%%%%%%%%%%%%%%%%%%%%%%%%%%%%%
\begin{frame}[fragile]\frametitle{Iteration Protocol}
The built-in function \lstinline|iter()| takes an iterable object and returns an iterator.
	\begin{lstlisting}
>>> x = iter([1, 2, 3])
>>> x
<list_iterator object at 0x1004ca850>
>>> next(x)
1
>>> next(x)
2
>>> next(x)
3
\end{lstlisting}
\end{frame}

%%%%%%%%%%%%%%%%%%%%%%%%%%%%%%%%%%%%%%%%%%%%%%%%%%%%%%%%%%%%%%%%%%%%%%%%%%%%%%%%%%%
\begin{frame}[fragile]\frametitle{Iteration Protocol}

\begin{lstlisting}
>>> next(x)
2
>>> next(x)
3
>>> next(x)
Traceback (most recent call last):
  File "<stdin>", line 1, in <module>
StopIteration
\end{lstlisting}
Each call to \lstinline|next()| returns the next element. When exhausted, it raises \lstinline|StopIteration|.
\end{frame}


%%%%%%%%%%%%%%%%%%%%%%%%%%%%%%%%%%%%%%%%%%%%%%%%%%%%%%%%%%%%%%%%%%%%%%%%%%%%%%%%%%%
\begin{frame}[fragile]\frametitle{Reading Lines}
Without Iterator:
\begin{lstlisting}
with open('data.txt', 'r') as f:
    while True:
        line = f.readline()
        if not line: 
            break
        else: 
            print(line.rstrip())
\end{lstlisting}

\end{frame}

%%%%%%%%%%%%%%%%%%%%%%%%%%%%%%%%%%%%%%%%%%%%%%%%%%%%%%%%%%%%%%%%%%%%%%%%%%%%%%%%%%%
\begin{frame}[fragile]\frametitle{Reading Lines}
With Iterator:
\begin{lstlisting}
with open('data.txt', 'r') as f:
    for line in f:
        print(line.rstrip())
\end{lstlisting}

\end{frame}


%%%%%%%%%%%%%%%%%%%%%%%%%%%%%%%%%%%%%%%%%%%%%%%%%%%%%%%%%%%%%%%%%%%%%%%%%%%%%%%%%%%
\begin{frame}[fragile]\frametitle{Iterator Class}
    \begin{itemize}
    \item Iterators are implemented as classes. 
    \item Here is an iterator that works like built-in range function.
	\begin{lstlisting}
class yrange:
    def __init__(self, n):
        self.i = 0
        self.n = n

    def __iter__(self):
        return self

    def __next__(self):
        if self.i < self.n:
            i = self.i
            self.i += 1
            return i
        else:
            raise StopIteration()
\end{lstlisting}
    \end{itemize}
\end{frame}

%%%%%%%%%%%%%%%%%%%%%%%%%%%%%%%%%%%%%%%%%%%%%%%%%%%%%%%%%%%%%%%%%%%%%%%%%%%%%%%%%%%
\begin{frame}[fragile]\frametitle{Iterator Class}
    \begin{itemize}
    \item \lstinline|__iter__| makes an object iterable; \lstinline|__next__| returns the next value.
    \item The built-in \lstinline|next()| function calls \lstinline|__next__()| on the iterator.
    \item When no more elements exist, \lstinline|__next__| must raise \lstinline|StopIteration|.
	\begin{lstlisting}
>>> y = yrange(3)
>>> next(y)
0
>>> next(y)
1
>>> next(y)
2
>>> next(y)
Traceback (most recent call last):
  File "<stdin>", line 1, in <module>
  File "<stdin>", line 14, in __next__
StopIteration
\end{lstlisting}
    \end{itemize}
\end{frame}

%%%%%%%%%%%%%%%%%%%%%%%%%%%%%%%%%%%%%%%%%%%%%%%%%%%%%%%%%%%%%%%%%%%%%%%%%%%%%%%%%%%
\begin{frame}[fragile]\frametitle{Iteration Usage}
Many built-in functions accept iterators as arguments.
\begin{lstlisting}
>>> list(yrange(5))
[0, 1, 2, 3, 4]
>>> sum(yrange(5))
10
\end{lstlisting}
In the above case, both the iterable and iterator are the same object. Notice that the \_\_iter\_\_ method returned self. It need not be the case always.
\end{frame}

%%%%%%%%%%%%%%%%%%%%%%%%%%%%%%%%%%%%%%%%%%%%%%%%%%%%%%%%%%%%%%%%%%%%%%%%%%%%%%%%%%%
\begin{frame}[fragile]\frametitle{Iteration Usage}
Iterable and Iterator can be different
\begin{lstlisting}
class zrange:
    def __init__(self, n):
        self.n = n

    def __iter__(self):
        return zrange_iter(self.n)

\end{lstlisting}
\end{frame}

%%%%%%%%%%%%%%%%%%%%%%%%%%%%%%%%%%%%%%%%%%%%%%%%%%%%%%%%%%%%%%%%%%%%%%%%%%%%%%%%%%%
\begin{frame}[fragile]\frametitle{Iteration Usage}
Iterable and Iterator can be different
\begin{lstlisting}
class zrange_iter:
    def __init__(self, n):
        self.i = 0
        self.n = n

    def __iter__(self):
        # Iterators are iterables too.
        # Adding this functions to make them so.
        return self

    def __next__(self):
        if self.i < self.n:
            i = self.i
            self.i += 1
            return i
        else:
            raise StopIteration()
\end{lstlisting}
\end{frame}


%%%%%%%%%%%%%%%%%%%%%%%%%%%%%%%%%%%%%%%%%%%%%%%%%%%%%%%%%%%%%%%%%%%%%%%%%%%%%%%%%%%
\begin{frame}[fragile]\frametitle{Iteration Usage}
If both iteratable and iterator are the same object, it is consumed in a single iteration.
\begin{lstlisting}
>>> y = yrange(5)
>>> list(y)
[0, 1, 2, 3, 4]
>>> list(y)
[]
>>> z = zrange(5)
>>> list(z)
[0, 1, 2, 3, 4]
>>> list(z)
[0, 1, 2, 3, 4]
\end{lstlisting}
\end{frame}



%%%%%%%%%%%%%%%%%%%%%%%%%%%%%%%%%%%%%%%%%%%%%%%%%%%%%%%%%%%%%%%%%%%%%%%%%%%%%%%%%%%
\begin{frame}[fragile]\frametitle{Context Managers (\texttt{with} statement)}
A \textbf{context manager} manages resource setup and teardown automatically --- even if an exception occurs.
\begin{lstlisting}
# Without context manager
f = open('data.txt')
try:
    data = f.read()
finally:
    f.close()

# With context manager (preferred)
with open('data.txt') as f:
    data = f.read()   # file is closed automatically on exit
\end{lstlisting}
Implement your own using \lstinline|__enter__| / \lstinline|__exit__| (or \lstinline|contextlib.contextmanager|):
\begin{lstlisting}
from contextlib import contextmanager

@contextmanager
def timer():
    import time; t = time.time()
    yield
    print(f"Elapsed: {time.time()-t:.3f}s")

with timer():
    sum(range(10_000_000))
\end{lstlisting}
\end{frame}

%%%%%%%%%%%%%%%%%%%%%%%%%%%%%%%%%%%%%%%%%%%%%%%%%%%%%%%%%%%%%%%%%%%%%%%%%%%%%%%%%%%
\begin{frame}[fragile]\frametitle{Iteration Advantages}
    \begin{itemize}
    \item Cleaner code
    \item     Iterators can work with infinite sequences
    \item     Iterators save resources
    \end{itemize}
   \end{frame}
   
%%%%%%%%%%%%%%%%%%%%%%%%%%%%%%%%%%%%%%%%%%%%%%%%%%%%%%%%%%%%%%%%%%%%%%%%%%%%%%%%%%%
\begin{frame}[fragile]\frametitle{Iteration Assignment}
Write an iterator class \lstinline|reverse_iter| that takes a list and iterates it from the reverse direction:
\begin{lstlisting}
>>> it = reverse_iter([1, 2, 3, 4])
>>> next(it)
4
>>> next(it)
3
>>> next(it)
2
>>> next(it)
1
>>> next(it)
Traceback (most recent call last):
  File "<stdin>", line 1, in <module>
StopIteration
\end{lstlisting}
\end{frame}





